\bookchapter{Chapter 34: In-Depth Guide to Car-Phone Connectivity — The Fusion of Phone and Automobile}

\chapterintro{Cars are evolving from “four wheels plus a sofa” to “four wheels plus a smartphone.” Whether you drive a gas car or an EV, the relationship between your phone and car infotainment system determines your driving experience. This chapter clarifies the three routes for car-phone connectivity, the characteristics of each camp's system, and tips that make driving easier.}

\booksection{I. Three Routes for Car-Phone Connectivity}

There are currently three approaches to connecting phones with cars. Understanding them means understanding the entire ecosystem:

\begin{itemize}
  \item \textbf{Phone Casting to Car Screen}: The car screen acts as just a "monitor" and "speaker" — computing and map data come from the phone. Representative solutions: Apple CarPlay, Google Android Auto, Huawei HiCar, Baidu CarLife. The advantage is that the phone ecosystem is directly brought to the car without reinstalling apps; the disadvantage is that the car's own hardware capabilities go unused.
  \item \textbf{Deep Integration of Phone and Car System}: The car system is independent but shares accounts, apps, and data with the phone, and can even "hand off" tasks. Representative solutions: HarmonyOS Cockpit (AITO/Luxeed), BYD DiLink, NIO OS, Li Auto system. The advantage is an integrated experience; the disadvantage is being locked into one brand ecosystem.
  \item \textbf{CarPlay 2.0 Deep Integration}: The next-generation CarPlay released by Apple in 2025 no longer just "casts" — it takes over core functions like the instrument cluster, climate control, and seat adjustment. Ferrari, Aston Martin, and Porsche are among the first to adopt it, marking a new era where "phone manufacturers drive the car infotainment system."
\end{itemize}

\booksection{II. CarPlay / Android Auto / HiCar Comparison}

The three mainstream casting solutions each have their strengths. Use this table when choosing a car or phone:

\begin{center}
\begin{tabularx}{\linewidth}{|X|X|X|X|X|}
\hline
\textbf{Solution} & \textbf{Supported Phone OS} & \textbf{Supported Car Models} & \textbf{Feature Richness} & \textbf{China Availability} \tabularnewline
\hline
Apple CarPlay & iPhone only & Very wide, most car models support it & High — maps/music/phone/messages all work & Good, works natively \tabularnewline
Android Auto & Android (domestic phones need service framework) & Fairly wide, mainly joint venture/imported cars & High, comparable to CarPlay & Average, requires VPN to activate \tabularnewline
Huawei HiCar & Mainly Huawei phones, some support other Android & Domestic cars and some joint venture models & Fairly high, can integrate with HarmonyOS Cockpit & Excellent, local ecosystem \tabularnewline
Baidu CarLife & Both iOS and Android & Fairly wide, common retrofit for older cars & Medium, updates are slow & Excellent, local ecosystem \tabularnewline
\hline
\end{tabularx}
\end{center}

\booksection{III. HarmonyOS Cockpit: Seamless Flow Between Phone and Car}

If CarPlay is “casting phone to car,” HarmonyOS Cockpit is “phone and car are already family.” The collaboration between Huawei phones and AITO/Luxeed car systems represents the highest level of domestic car-phone connectivity:

\begin{itemize}
  \item \textbf{Super Desktop}: The car system directly accesses apps on the phone — no need to reinstall on the car. Music playing on the phone and navigation routes automatically continue on the car screen when you get in, and return to the phone when you get out — seamless throughout.
  \item \textbf{App Handoff}: Navigation set on the phone is automatically taken over by the car system when you close the door; when you arrive at the parking lot, the phone continues walking navigation to the final destination. The "last mile" is no longer a gap.
  \item \textbf{Shared Account Between Phone and Car}: Log into your Huawei account once, and the car system automatically syncs your music membership, map favorites, schedule, and contacts. The car and phone are different terminals of the same identity, not two independent systems.
\end{itemize}

\booksection{IV. Car Systems from New Automakers}

Domestic EV infotainment systems have evolved from “navigation + music” to mobile smart terminals, each with distinctive approaches:

\begin{itemize}
  \item \textbf{NIO OS}: The NOMI voice assistant is an industry benchmark — it can recognize multiple commands and contextual continuous dialogue, with emotional expressions. The NIO Phone turns the phone into a "car key + second screen," deeply integrated with the NIO ecosystem.
  \item \textbf{Li Auto System}: Multi-screen interaction is the highlight — the front passenger screen can independently play shows, and the rear screen can play cartoons. "Task Master" lets you customize scenarios, like "auto-pause music when door opens, resume when closed."
  \item \textbf{Xiaomi HyperOS Car System}: Deeply integrated with Xiaomi phones, watches, and home appliances. Alarms set on the phone remind you on the car system; when navigating home, the home AC turns on in advance. HyperOS connects the "person-car-home ecosystem."
  \item \textbf{BYD DiLink}: The signature rotating screen — vertical for short videos, horizontal for maps, adapting to app orientation. Built-in app store leads in ecosystem richness among traditional automakers.
\end{itemize}

\booksection{V. Phone as Car System — Can It Work?}

If your car system is too old or too hard to use, can you just use your phone instead? The answer: yes, but it depends on the scenario.

\begin{itemize}
  \item \textbf{Motorola Ready For / Desktop Mode}: Some Motorola, Samsung, and Huawei phones support "Desktop Mode" — connect a monitor and it becomes a simple computer. In theory, it could connect to a car screen, but real-world compatibility varies, and it lacks car-specific interactions (like steering wheel buttons).
  \item \textbf{Phone Mount vs. Built-in Car System}: If the built-in car system lags or maps are outdated, a sturdy phone mount + phone navigation is a practical choice. But note: (1) small phone screen is hard to see while driving; (2) summer heat can cause the phone to thermal-throttle or shut down; (3) long navigation drains battery fast — you need a car charger.
  \item \textbf{Recommendation}: For short trips or temporary use, a phone mount is fine; for daily commuting and long trips, upgrading to a CarPlay/HiCar-enabled system or installing an aftermarket unit is better for both experience and safety.
\end{itemize}

\begin{tcbwarn}{Security Reminder}
Regardless of which solution you use, operating your phone while driving violates traffic rules. A phone mount is only for “viewing navigation” — not for “operating the phone.” Set destinations, switch songs, etc. only after pulling over, or use voice control.

\end{tcbwarn}

\booksection{VI. Practical Tips}

\begin{itemize}
  \item \textbf{Unstable Wireless CarPlay Connection}: (1) Turn off car Bluetooth and re-pair; (2) Delete paired device records on the car side and reconnect; (3) Check if your phone is connected to multiple Bluetooth devices causing conflicts; (4) Update car firmware and phone system to the latest version; (5) If all else fails, switch to wired CarPlay — more stable.
  \item \textbf{Car Data Plan Used Up}: (1) Turn on phone hotspot and connect the car system to it; (2) In car settings, switch "high data consumption" apps (like video and streaming music) to Wi-Fi; (3) Regularly clear the car's offline map cache.
  \item \textbf{Driving Mode Auto-Trigger}: (1) iOS "Focus" set to "auto-enable when connected to car Bluetooth," automatically blocking non-urgent notifications; (2) Android "Driving Mode" or third-party apps (like Tasker) can auto-mute and auto-reply to messages when connected to Bluetooth; (3) Huawei phones can set "smart payment" to auto-switch music to the car system when you get in.
\end{itemize}

\begin{tcbinfo}{Tips When Choosing a Car}
If you care about the infotainment experience, be sure to test at the dealership: (1) Bring your phone and test CarPlay/HiCar connection stability; (2) Test voice assistant response speed and command recognition; (3) Check if the app store has your frequently used apps; (4) Pay attention to OTA update frequency — good car systems are built through software iteration.

\end{tcbinfo}

\begin{tcbpractice}{Hands-On: Optimize Your Car Connectivity Experience}
\begin{enumerate}
  \item Confirm which connectivity solution your phone and car support (CarPlay / HiCar / Android Auto / CarLife).
  \item Clear the car's Bluetooth pairing list once, keeping only your commonly used phones to reduce connection conflicts.
  \item Set up a "Driving Focus" that auto-activates when connected to car Bluetooth.
  \item Download the car brand app (e.g., HarmonyOS Smart Driving, Li Auto, NIO App) to see if you can remotely control AC and check vehicle status.
\end{enumerate}

\end{tcbpractice}

\begin{tcbsummary}{Chapter Summary}
\begin{itemize}
  \item Three connectivity routes: casting (CarPlay etc.), deep integration (HarmonyOS Cockpit etc.), CarPlay 2.0 deep integration
  \item CarPlay has the widest ecosystem; HiCar has the best domestic experience; Android Auto requires workarounds in China; CarLife suits older cars
  \item HarmonyOS Cockpit's Super Desktop, App Handoff, and shared accounts represent the deep integration direction
  \item NIO NOMI, Li Auto multi-screen, Xiaomi person-car-home, BYD rotating screen — each has its own style
  \item Using a phone as car system is only suitable for temporary scenarios — upgrade or install aftermarket for long-term use
  \item Practical tips: CarPlay reconnection, phone hotspot for data, Bluetooth-triggered driving mode
\end{itemize}

\end{tcbsummary}
